\chapter{导论}

\section{研究背景}

N皇后问题是组合数学中最古老、最经典的问题之一。1848年，德国国际象棋爱好者
Bezzel在《柏林国际象棋杂志》上首次提出了著名的八皇后问题\cite{Bezzel1848}：
在$8\times8$的国际象棋棋盘上放置8个皇后，使得任意两个皇后互不攻击——即不共享
同一行、同一列或同一对角线。这个问题迅速吸引了数学界的关注，包括伟大的数学家高斯
（Gauss）在内都参与了讨论。高斯与Nauck等人很快算出了$8\times8$棋盘上全部92个
解。此后，N皇后问题——将上述问题推广到$N\times N$棋盘，放置$N$个互不攻击的皇后
——成为组合枚举理论的经典案例，历经一个半世纪而不衰。

从数学结构上看，N皇后问题是一个典型的约束满足问题（Constraint Satisfaction
Problem, CSP）：需要在$N^{2}$个离散格点上放置$N$个占据粒子，同时满足行、列、
两族对角线共四组``互不共存''约束。设$Q(N)$为N皇后问题解的总数。对小尺寸$N$，
$Q(N)$可以通过精确回溯搜索枚举：$Q(1)=1$，$Q(2)=Q(3)=0$，$Q(4)=2$，
$Q(5)=10$，$Q(6)=4$，$Q(7)=40$，$Q(8)=92$，$Q(9)=352$，$Q(10)=724$等。
然而，$Q(N)$随$N$的增加呈超指数增长，精确枚举很快陷入组合爆炸。到2025年，
利用GPU大规模并行回溯搜索，$Q(N)$的精确枚举纪录仅被推至$N=27$\cite{Yao2025}，
凸显了直接枚举方法固有的指数复杂度瓶颈。

N皇后问题的一个核心理论问题是确定$Q(N)$的渐近行为——当$N$趋于无穷大时，
$Q(N)$的增长率是多少？这个问题在数学上极为困难，经过一个半世纪的探索，
直到2022年才由以色列数学家Simkin给出了严格证明\cite{Simkin2022}。
Simkin的工作分两步：首先引入``queenons''——定义在单位正方形上的概率测度
——作为$N\to\infty$极限下的皇后占据密度场，将问题转化为这些极限对象空间上的
凸优化问题；然后利用熵方法（随机抽样+联合熵界）建立上界，
再通过精巧的随机化构造算法建立下界。Simkin的主要结果是：
%
\begin{equation}
  Q(N)\sim \left(\frac{N}{e^{\gamma}}\right)^{N},
  \label{eq:QN}
\end{equation}
%
其中常数$\gamma$位于区间$[1.939,1.945]$内。这一常数被后人称为Simkin常数，
是N皇后问题的核心组合学参数。

在Simkin的证明之前，Bowtell与Keevash\cite{Bowtell2023}首先建立了$Q(N)$的
指数下界，Luria与Simkin\cite{Luria2021}进一步给出更精细的估计，
这些前期突破为最终的严格证明铺平了道路。随后，Nobel等人\cite{Nobel2023}
于2023年利用大规模Newton方法数值求解Simkin的凸优化程序，
将$\gamma$的界收紧到惊人的精度：$\gamma=1.94400(1)$，精确度约为$10^{-7}$量级，
为后续的数值验证提供了高精度的参照基准。

Simkin的渐近公式$Q(N)\sim(N/e^{\gamma})^{N}$具有一个极具启发性的形式。
取对数后，每皇后的基态熵为
$s_{0}=\lim_{N\to\infty}\ln Q(N)/N=\ln N-\gamma$。
这一表达式的结构值得深入分析：在每行放置一个皇后的前提下（行使皇后数恰好为1的
行约束由构造自然满足），每个皇后可以从$N$列中独立选择一列，给出$\ln N$单位的熵。
N皇后问题的熵相对于这一自由放置值恰好减少了常数$\gamma$。换句话说，
约束的累积效应可被精确量化为每皇后约1.944单位的熵损失。

进一步的分析可以揭示约束的层级结构\cite{Knuth2022}：
如果仅施加列不可重复约束（每列恰好一个皇后），配置空间为所有$N!$个列排列
（即著名的``车问题''或rook problem）。由Stirling公式
$\ln N!=N\ln N-N+O(\ln N)$，在$N\to\infty$极限下每皇后的熵为
$s_{\text{perm}}=\ln N-1+O(\ln N/N)$。
因此，列约束的熵代价在大$N$下趋于恰好1个单位——这就是著名的Stirling修正。
两族对角线的进一步引入使熵减少至$\ln N-\gamma$，对应额外代价
$\gamma-1\approx0.94$单位。

这种每皇后熵按约束层级逐级减少一个$N$无关常数的模式，强烈暗示N皇后问题
背后存在一个简洁的统计力学结构。从统计力学角度看，N皇后问题的解是零能量基态
的简并集合，简并度即为$Q(N)$。引入温度后，系统可以在攻击构型（非零能量）
与非攻击构型（零能量）之间热激发。有限温度下的热力学行为——包括能量、比热、
熵等基本热力学量的温度依赖关系——构成了反映约束物理本质的丰富图景。
通过Monte Carlo模拟在有限温度下探测这些热力学量，并利用热力学关系
$S(\infty)-S(0)=\int_{0}^{\infty}C_{v}/T\,\mathrm{d}T$将高温与基态联系
起来，便有可能从纯热力学数据中反推出基态简并度——从而为Simkin常数$\gamma$
提供一条独立于组合枚举的热力学测定路径。这正是本文的核心研究思路。

\section{国内外研究现状}

\subsection{N皇后问题的组合枚举与渐近分析}

N皇后问题的研究可追溯至19世纪中叶。Knuth在《计算机程序设计艺术》第4B卷
\cite{Knuth2022}中系统地总结了N皇后问题的回溯算法、对称性约简以及与拉丁方阵、
精确覆盖等其他经典组合问题的内在联系。在精确枚举方面，从20世纪60年代起，
随着计算机技术的发展，$Q(N)$的已知范围不断扩展。最新的里程碑是Yao和Li
\cite{Yao2025}于2025年利用GPU大规模并行回溯搜索，将$Q(N)$的精确枚举推至$N=27$，
这是当前已知的世界纪录。

在渐近分析方面，$Q(N)$的渐近行为长期是组合数学的悬案。突破始于2021年：
Bowtell和Keevash\cite{Bowtell2023}首次建立了$Q(N)$的指数下界的严格证明，
Luria和Simkin\cite{Luria2021}进一步给出更精细的估计。Simkin\cite{Simkin2022}
的证明方法是一套精巧的组合：首先定义queenons——单位正方形$[0,1]^{2}$上的
概率测度；然后将$Q(N)$的渐近估计转化为queenons空间上的凸优化问题；
上界利用熵方法，下界则通过随机化构造算法显式生成大量有效的皇后配置。
Nobel等人\cite{Nobel2023}随后将Simkin的凸优化问题表述为大规模的非线性规划，
并利用不可行起始Newton方法结合对数障碍函数进行数值求解，
将$\gamma$的界收紧到$\gamma=1.94400(1)$。

\subsection{N皇后问题的统计物理学研究}

将N皇后问题映射为统计力学模型并进行Monte Carlo模拟，是一条相对独立的研究路线。
Zhang和Ma\cite{Zhang2009}于2009年发表了具有里程碑意义的工作：他们将N皇后问题
视为正则系综中的格点气体，利用Monte Carlo模拟结合逐步偏置抽样估计$Q(N)$。
Polson和Sokolov\cite{Polson2024}于2024年延续并深化了这一研究方向。他们明确提出了
将N皇后问题视为格点气体的视角，并推导了高温极限下的若干解析结果。特别地，
他们利用攻击线的线性分解和期望的线性性，导出了高温下每皇后平均能量的精确表达式。

\subsection{蒙特卡洛方法与张量网络方法}

在数值方法层面，本论文依赖于两个成熟的技术体系。其一是Metropolis等人
\cite{Metropolis1953}于1953年提出的Monte Carlo算法，它是统计物理模拟的基石。
对于固定粒子数的格点气体系统，Kawasaki\cite{Kawasaki1966}于1966年提出的
粒子-空位交换动力学（Kawasaki dynamics）是天然的选择。在统计误差估计方面，
Efron\cite{Efron1982}提出的jackknife重采样方法为相关时间序列提供了可靠的
误差估计方案。Binder和Heermann的经典著作\cite{Binder2010}对Monte Carlo模拟
的物理基础做了系统的阐述。

其二是张量网络（Tensor Network）方法。起源于量子多体物理中的密度矩阵重正化群
（DMRG），矩阵乘积算符（Matrix Product Operator, MPO）提供了一种编码
一维序列约束的紧凑表达方式\cite{Frowis2010}。Kourtis等人\cite{Kourtis2019}
在2019年的工作中展示了如何利用张量网络对约束满足问题进行精确计数。
对于N皇后问题，每个格点同时受到来自四个方向的约束，产生一个rank-9的局部张量
——恰好具有17个非零元素。通过逐行收缩此张量网络\cite{Xiang2024}，
可以精确地计算出$Q(N)$，为Monte Carlo模拟提供无噪声的精确基准。

\section{本文的研究目的与结构安排}

本论文的核心目标是：建立N皇后问题的统计力学框架，通过大规模Monte Carlo模拟
绘制其完整的热力学相图，并利用热力学积分方法从纯模拟数据中独立提取Simkin常数
$\gamma$，从而为该基本组合常数提供一个独立于组合数学的热力学测定途径。

本文的主要创新点包括：（1）首次通过系统的热力学积分从Monte Carlo数据中独立
反推出Simkin常数$\gamma$，在$N=1024$处达到$0.11\%$的高精度；（2）系统揭示了
N皇后格点气体的约束层级结构，展示了各几何约束如何贡献可精确量化的每皇后熵代价；
（3）给出了高温极限$E/N\to5/3$的严格有限$N$表达式；（4）建立了张量网络
精确枚举方案，揭示了N皇后转移矩阵的幂零性质。

本文的结构安排如下：第二章详细介绍格点气体模型的构建和Monte Carlo模拟方法。
第三章给出解析结果，包括基态熵的约束层级分析和高温极限的严格推导。
第四章展示有限温度Monte Carlo模拟结果。
第五章进行热力学积分，从纯Monte Carlo数据中独立提取Simkin常数。
第六章建立张量网络精确枚举方案，揭示N皇后转移矩阵的幂零性及其方法论意义。
第七章总结全文并讨论未来展望。

本研究的主要结果已与合作者Hai-Jun Liao和Lei Wang共同撰写为英文预印本
\textit{Statistical Mechanics of the N-Queens Problem}，于2026年5月发布在
arXiv\cite{Liu2026arxiv}上。本文在内容组织与论述细节上对该工作进行了
独立的中文重述与扩展，但核心结论与数据一致。
